% Generated draft: requires mathematical and source review.
\documentclass[11pt]{article}
\usepackage[margin=25mm]{geometry}
\usepackage{fontspec}
\setmainfont{Latin Modern Roman}
\usepackage{amsmath,amssymb,mathtools,unicode-math}
\setmathfont{Latin Modern Math}
\usepackage{enumitem,needspace,xurl,hyperref}
\setlength{\parindent}{0pt}
\setlength{\parskip}{5pt}
\setlength{\emergencystretch}{3em}
\newcommand{\N}{\mathbb{N}}
\newcommand{\Z}{\mathbb{Z}}
\newcommand{\Q}{\mathbb{Q}}
\newcommand{\R}{\mathbb{R}}
\newcommand{\C}{\mathbb{C}}
\newcommand{\F}{\mathbb{F}}
\newcommand{\A}{\mathbb{A}}
\newcommand{\PP}{\mathbb{P}}
\newcommand{\E}{\mathbb{E}}
\newcommand{\eps}{\varepsilon}
\newcommand{\dd}{\,\mathrm{d}}
\begin{document}
\section*{UnsolvedMath ID: 30001471}
\textbf{Source problem code:} catalogue entry
\subsection{Definitions and mathematical statement}
A global field is a finite extension of either \(\Q\) or \(\F_q(t)\), where \(\F_q\) is a finite field. Let \(K\) be a global field, let \(S\) be a finite set of inequivalent valuations of \(K\), and let \(G\) be a noncommutative absolutely almost simple \(K\)-isotropic algebraic group. Here \(K\)-isotropic means that \(G\) contains a nontrivial \(K\)-split torus. Write
\[
\Gamma=G(\mathcal O_S)\subset G_S,
\]
and retain the source's hypothesis that \(\Gamma\) is an arithmetic lattice in \(G_S\).

\textbf{Statement requested by the source.} For every input scale \(r_1\), do there exist a linear polynomial \(f\) and an output scale \(r_2\) such that, for every dimension \(n<k(G,S)\) and every coarse \(n\)-manifold \(\Sigma\subset G_S\) of scale \(r_1\) satisfying \(\partial\Sigma\subset\Gamma\), there is a coarse \(n\)-manifold \(\Sigma'\subset\Gamma\) of scale \(r_2\), of the same topological type as \(\Sigma\), satisfying
\[
\partial\Sigma'=\partial\Sigma,
\qquad
\operatorname{vol}(\Sigma')\leq f\bigl(\operatorname{vol}(\Sigma)\bigr)?
\]
The displayed quantifier order requires \(f\) and \(r_2\) to be independent of \(\Sigma\); it reads the dimension restriction as applying to all admissible \(n\) after these choices.

\textbf{Unresolved qualifications.} The supplied record does not define \(\mathcal O_S\), \(G_S\), or \(k(G,S)\), or specify the metric, the meaning of a coarse manifold and its scale, or the volume convention. It also does not specify whether equality of boundaries means equality of parametrized boundary data or equality of their images, nor whether the topological identification must respect that boundary. These source-dependent definitions must be supplied before the question becomes a fully specified mathematical proposition. In particular, no formula for \(k(G,S)\) is asserted here. No additional restrictions on \(S\), such as inclusion of archimedean places, are imposed by this draft; the asserted lattice hypothesis is retained explicitly.
\subsection{Short English statement}
The question asks whether every coarse manifold in the ambient group, below the unspecified threshold \(k(G,S)\) and with boundary in the arithmetic lattice, can be replaced inside that lattice with the same boundary and topological type, using a uniformly controlled scale and a linear volume bound. Several essential geometric conventions are absent from the supplied record and require review.
\subsection{Sources}
\begin{itemize}
\item Local UnsolvedMath registry, supplied record with permanent ID \(30001471\), source problem code \texttt{catalogue entry}, entitled \emph{Coarse-Manifold Fillings in S-Arithmetic Lattices}, in the Oberwolfach Reports Open Problems collection. The record gives proposed year 2010 and names no individual proposer.
\item Supplied report citation: \emph{Non-positive Curvature and Geometric Structures in Group Theory} (2010), Oberwolfach Reports, pp.\ 1165--1224, DOI \texttt{10.4171/owr/2010/20}. Source URL and supplied literature-source link: \url{https://doi.org/10.4171/owr/2010/20}.
\item Supplied literature-source attribution: Bestvina--Eskin--Wortman, \url{https://doi.org/10.4171/JEMS/419}. The supplied record does not give a paper title or more detailed bibliographic information.
\item Literature assessment dated \texttt{2026-08-21}: the supplied curation labels the problem open and states that Bestvina--Eskin--Wortman obtain partial results toward the coarse-manifold filling conjecture while the general statement remains unproved. This is supplied dated curation, not an independently verified literature assessment; its time-sensitive status and cited claims require independent verification.
\end{itemize}
\end{document}
